\documentclass[../../main.tex]{subfiles}

\begin{document}

\section{紧群上的不变 Borel 测度: von Neumann 定理}\label{section:ch22-haar}

紧拓扑空间 $X$ 上的 Borel 测度是 $\mathcal{B}(X)$ 上的有限测度, $\mathcal{B}(X)$ 即为包含 $X$ 上的拓扑的最小 $\sigma$ 代数. 我们现在考虑紧群上的 Borel 测度以及它们与群运算的关系.

\begin{lemma}[左平移测度]\label{lemma:ch22-left-translate}
令 $\mathcal{G}$ 为紧群而 $\mu$ 是 $\mathcal{B}(\mathcal{G})$ 上的 Borel 测度. 对 $g\in\mathcal{G}$, 定义集函数 $\mu_{g}:\mathcal{B}(\mathcal{G})\to[0,\infty)$ 为
\begin{align*}
\text{对所有 }A\in\mathcal{B}(\mathcal{G}),\quad\mu_{g}(A)=\mu(g\cdot A).
\end{align*}
则 $\mu_{g}$ 是 Borel 测度. 若 $\mu$ 是 Radon 的, $\mu_{g}$ 也是. 此外, 若 $\pi$ 是 $\mathcal{G}$ 在 $C(\mathcal{G})$ 上的正则表示, 则
\begin{align}
\text{对所有 }f\in C(\mathcal{G}),\quad\int_{\mathcal{G}}\pi(g)f\,d\mu=\int_{\mathcal{G}}f\,d\mu_{g}.
\label{align:ch22-translate-integral}
\end{align}
\end{lemma}

\begin{proof}
令 $g$ 属于 $\mathcal{G}$. 观察到用 $g$ 做左乘定义了 $\mathcal{G}$ 映上 $\mathcal{G}$ 的拓扑同胚. 由此我们推出 $A$ 是 Borel 集当且仅当 $g\cdot A$ 是 Borel 集. 因此集函数 $\mu_{g}$ 在 $\mathcal{B}(\mathcal{G})$ 上是恰当定义的. 显然, $\mu_{g}$ 继承了 $\mu$ 的可数可加性, 由于 $\mu_{g}(\mathcal{G})=\mu(\mathcal{G})<\infty$, $\mu_{g}$ 是 Borel 测度. 现在假定 $\mu$ 是 Radon 测度. 为证明 $\mu_{g}$ 的内正则性, 令 $\mathcal{O}$ 为 $\mathcal{G}$ 中的开集而 $\varepsilon>0$. 由于 $\mu$ 是内正则的, $g\cdot\mathcal{O}$ 是开的, 存在包含于 $g\cdot\mathcal{O}$ 的紧集 $K$ 使得 $\mu(g\cdot\mathcal{O}\sim K)<\varepsilon$. 因此 $K'=g^{-1}\cdot K$ 是紧的, 包含于 $\mathcal{O}$ 且 $\mu_{g}(\mathcal{O}\sim K')<\varepsilon$. 因此 $\mu_{g}$ 是内正则的. 用类似的方法可证明 $\mu_{g}$ 是外正则的. 因此 $\mu_{g}$ 是 Radon 测度.

我们现在证明\eqref{align:ch22-translate-integral}. 积分是线性的. 因此, 若\eqref{align:ch22-translate-integral}对 Borel 集的特征函数成立, 它对简单 Borel 函数也成立. 我们从简单逼近定理和有界收敛定理推出, 若\eqref{align:ch22-translate-integral}对简单 Borel 函数成立, 则它对所有 $f\in C(\mathcal{G})$成立. 因此仅须对 Borel 集 $A$ 的特征函数 $f=\chi_{A}$ 的情形证明\eqref{align:ch22-translate-integral}. 然而, 对这样的函数,
\begin{align*}
\int_{\mathcal{G}}\pi(g)f\,d\mu=\mu(g\cdot A)=\int_{\mathcal{G}}f\,d\mu_{g}.
\end{align*}
\end{proof}

\begin{definition}[左不变测度与概率测度]\label{definition:ch22-left-invariant}
令 $\mathcal{G}$ 为紧群. Borel 测度 $\mu:\mathcal{B}(\mathcal{G})\to[0,\infty)$ 称为是左不变的, 若
\begin{align}
\text{对所有 }g\in\mathcal{G}\text{ 与 }A\in\mathcal{B}(\mathcal{G}),\quad\mu(A)=\mu(g\cdot A).
\label{align:ch22-left-invariant-def}
\end{align}
它称为是概率测度, 若 $\mu(\mathcal{G})=1$.
\end{definition}

右不变测度可类似定义. 若考虑 $\mathbf{R}^{n}$ 为加法运算的拓扑群, 我们证明 $\mathbf{R}^{n}$ 上的 Lebesgue 测度限制在 $\mathcal{B}(\mathbf{R}^{n})$ 上关于加法是左不变的, 即对 $\mathbf{R}^{n}$ 的每个 Borel 子集 $E$ 与每个点 $x\in\mathbf{R}^{n}$, $\mu_{n}(E+x)=\mu_{n}(E)$. 当然, 这对 $\mathbf{R}^{n}$ 的任何 Lebesgue 可测子集 $E$ 也成立.

\begin{proposition}[不变 Radon 概率测度的存在性]\label{proposition:ch22-existence}
在每个紧群 $\mathcal{G}$ 上, 存在 $\mathcal{B}(\mathcal{G})$ 上的左不变 Radon 概率测度, 也存在右不变的 Radon 概率测度.
\end{proposition}

\begin{proof}
定理\cref{theorem:ch22-kakutani}告诉我们存在概率泛函 $\psi\in[C(\mathcal{G})]^{*}$, 它在 $\mathcal{G}$ 在 $C(\mathcal{G})$ 上的正则表示的伴随下是不动的. 这意味着 $\psi(1)=1$ 且
\begin{align}
\text{对所有 }f\in C(\mathcal{G})\text{ 与 }g\in\mathcal{G},\quad\psi(f)=\psi(\pi(g^{-1})f).
\label{align:ch22-invariant-functional}
\end{align}
另一方面, 根据 Riesz-Markov 定理, $\mathcal{B}(\mathcal{G})$ 上存在唯一的 Radon 测度 $\mu$, 在
\begin{align}
\text{对所有 }f\in C(\mathcal{G}),\quad\psi(f)=\int_{\mathcal{G}}f\,d\mu
\label{align:ch22-riesz-markov}
\end{align}
的意义下表示 $\psi$.

因此, 根据\eqref{align:ch22-invariant-functional}, 对所有 $f\in C(\mathcal{G})$和 $g\in\mathcal{G}$,
\begin{align}
\psi(f)=\psi(\pi(g^{-1})f)=\int_{\mathcal{G}}\pi(g^{-1})f\,d\mu.
\label{align:ch22-integral-invariant}
\end{align}
因此, 根据引理\cref{lemma:ch22-left-translate}, 对所有 $f\in C(\mathcal{G})$和 $g\in\mathcal{G}$,
\begin{align*}
\psi(f)=\int_{\mathcal{G}}f\,d\mu_{g^{-1}}.
\end{align*}
同样由引理\cref{lemma:ch22-left-translate}, $\mu_{g^{-1}}$ 是 Radon 测度. 我们从泛函 $\psi$ 表示的唯一性推出,
\begin{align*}
\text{对所有 }g\in\mathcal{G},\quad\mu=\mu_{g^{-1}}.
\end{align*}
因此 $\mu$ 是左不变 Radon 测度. 因为 $\psi$ 是概率泛函, $\mu$ 是概率测度, 所以,
\begin{align*}
1=\psi(1)=\int_{\mathcal{G}}d\mu=\mu(\mathcal{G}).
\end{align*}
对偶方法(见习题\cref{exercise:ch22-25})证明了右不变 Radon 概率测度的存在性.
\end{proof}

\begin{definition}[Haar 测度]\label{definition:ch22-haar}
令 $\mathcal{G}$ 为拓扑群. $\mathcal{B}(\mathcal{G})$ 上的 Radon 测度称为是 Haar 测度, 若它是左不变的概率测度.
\end{definition}

\begin{theorem}[Von Neumann: Haar 测度的唯一性]\label{theorem:ch22-von-neumann}
令 $\mathcal{G}$ 为紧群. 则 $\mathcal{B}(\mathcal{G})$ 上存在唯一的 Haar 测度 $\mu$. 测度 $\mu$ 也是右不变的.
\end{theorem}

\begin{proof}
根据前一个命题, $\mathcal{B}(\mathcal{G})$ 上存在左不变的 Radon 概率测度 $\mu$ 和右不变的 Radon 概率测度 $\nu$. 我们宣称:
\begin{align}
\text{对所有 }f\in C(\mathcal{G}),\quad\int_{\mathcal{G}}f\,d\mu=\int_{\mathcal{G}}f\,d\nu.
\label{align:ch22-uniqueness}
\end{align}
一旦这得到证实, 我们从由对 Radon 测度的积分给出的 $C(\mathcal{G})$ 上有界线性泛函的表示的唯一性推出 $\mu=\nu$. 因此每个左不变的 Radon 测度等于 $\nu$. 因此仅有一个左不变的 Radon 测度且它是右不变的.

为证明\eqref{align:ch22-uniqueness}, 令 $f$ 属于 $C(\mathcal{G})$. 对 $(x,y)\in\mathcal{G}\times\mathcal{G}$ 定义 $h:\mathcal{G}\times\mathcal{G}\to\mathbf{R}$ 为 $h(x,y)=f(x\cdot y)$. 则 $h$ 是 $\mathcal{G}\times\mathcal{G}$ 上的连续函数. 此外, 乘积测度 $\mu\times\nu$ 定义在 $\mathcal{G}\times\mathcal{G}$ 的包含 $\mathcal{B}(\mathcal{G}\times\mathcal{G})$ 的子集的 $\sigma$ 代数上. 因此, 由于 $h$ 是可测的且在有限 $\mu\times\nu$ 测度集 $\mathcal{G}\times\mathcal{G}$ 上有界, 它关于 $\mathcal{G}\times\mathcal{G}$ 上的乘积测度 $\mu\times\nu$ 是可积的. 为证明\eqref{align:ch22-uniqueness}, 仅须证明:
\begin{align}
\int_{\mathcal{G}\times\mathcal{G}}h\,d[\nu\times\mu]=\int_{\mathcal{G}}f\,d\mu\text{ 与 }\int_{\mathcal{G}\times\mathcal{G}}h\,d[\mu\times\nu]=\int_{\mathcal{G}}f\,d\nu.
\label{align:ch22-fubini-pair}
\end{align}
然而, 由 Fubini 定理,
\begin{align*}
\int_{\mathcal{G}\times\mathcal{G}}h\,d[\nu\times\mu]=\int_{\mathcal{G}}\left[\int_{\mathcal{G}}h(x,\cdot)\,d\mu(y)\right]d\nu(x).
\end{align*}
根据 $\mu$ 的左不变性和\eqref{align:ch22-translate-integral},
\begin{align*}
\text{对所有 }x\in\mathcal{G},\quad\int_{\mathcal{G}}h(x,\cdot)\,d\mu(y)=\int_{\mathcal{G}}f\,d\mu.
\end{align*}
因此, 由于 $\nu(\mathcal{G})=1$,
\begin{align*}
\int_{\mathcal{G}\times\mathcal{G}}h\,d[\mu\times\nu]=\int_{\mathcal{G}}f\,d\mu\cdot\nu(\mathcal{G})=\int_{\mathcal{G}}f\,d\mu.
\end{align*}
类似的方法可证明\eqref{align:ch22-fubini-pair}右边的等式, 因此完成了定理的证明.
\end{proof}

我们这里的方法可被推广用于证明任何局部紧群 $\mathcal{G}$ 上存在左不变的 Haar 测度, 虽然它可以不是右不变的. 这里我们研究拓扑群上的拓扑, 确定它的测度论性质. 当然, 研究测度对拓扑的影响也是有趣的. 这一轮有趣的思想的进一步研究, 读 John von Neumann 的经典讲义《Invariant Measures》[\cite{vN91}]是有益的.

\begin{exercise}\label{exercise:ch22-19}
令 $\mu$ 为紧群 $\mathcal{G}$ 上的 Borel 概率测度. 证明 $\mu$ 是 Haar 测度当且仅当
\begin{align*}
\text{对所有 }g\in\mathcal{G},\ f\in C(\mathcal{G}),\quad\int_{\mathcal{G}}f\circ\varphi_{g}\,d\mu=\int_{\mathcal{G}}f\,d\mu,
\end{align*}
其中对所有 $g'\in\mathcal{G}$, $\varphi_{g}(g')=g\cdot g'$.
\end{exercise}

\begin{exercise}\label{exercise:ch22-20}
令 $\mu$ 为紧群 $\mathcal{G}$ 上的 Haar 测度. 证明 $\mu\times\mu$ 是 $\mathcal{G}\times\mathcal{G}$ 上的 Haar 测度.
\end{exercise}

\begin{exercise}\label{exercise:ch22-21}
令 $\mathcal{G}$ 为紧群, 其拓扑由度量给出. 证明存在 $\mathcal{G}$ 不变度量. (提示: 利用前两个习题且平均 $\mathcal{G}\times\mathcal{G}$ 上的度量.)
\end{exercise}

\begin{exercise}\label{exercise:ch22-22}
令 $\mu$ 为紧群 $\mathcal{G}$ 上的 Haar 测度. 若 $\mathcal{G}$ 有无穷多成员, 证明对每个 $g\in\mathcal{G}$, $\mu(\{g\})=0$. 若 $\mathcal{G}$ 是有限的, 直接描述 $\mu$.
\end{exercise}

\begin{exercise}\label{exercise:ch22-23}
证明: 若 $\mu$ 是紧群上的 Haar 测度, 则对 $\mathcal{G}$ 的每个开子集 $\mathcal{O}$, $\mu(\mathcal{O})>0$.
\end{exercise}

\begin{exercise}\label{exercise:ch22-24}
令 $S^{1}=\{z=e^{i\theta}\mid\theta\in\mathbf{R}\}$ 为具有复数乘法的群运算与承袭自欧氏圆盘的拓扑的圆周.
\begin{enumerate}[(i)]
\item 证明 $S^{1}$ 是拓扑群.
\item 定义 $\Lambda=\{(\alpha,\beta)\mid\alpha,\beta\in\mathbf{R},\ 0<\beta-\alpha<2\pi\}$. 对 $\lambda=(\alpha,\beta)\in\Lambda$, 定义 $I_{\alpha}=\{e^{i\theta}\mid\alpha<\theta<\beta\}$. 证明 $S^{1}$ 的每个真开子集是形如 $I_{\lambda}(\lambda\in\Lambda)$的集合的可数不交并.
\item 对 $\lambda=(\alpha,\beta)\in\Lambda$, 定义 $\mu(I_{\lambda})=(\beta-\alpha)/2\pi$. 定义 $\mu(S^{1})=1$. 用(ii)部分将 $\mu$ 延拓为定义在 $S^{1}$ 的开拓扑 $\mathcal{T}$ 上的集函数. 接着用前一章的命题 $9$ 证明 $\mu$ 可被延拓为 $\mathcal{B}(S^{1})$ 上的 Borel 测度.
\item 证明(ii)部分定义的测度是 $S^{1}$ 上的 Haar 测度.
\item 圆环 $T^{n}$ 是由具有乘积拓扑与群结构的 $S^{1}$ 的 $n$ 重笛卡儿积组成的拓扑群. $T^{n}$ 上的 Haar 测度是什么?
\end{enumerate}
\end{exercise}

\begin{exercise}\label{exercise:ch22-25}
令 $\mu$ 为拓扑群 $\mathcal{G}$ 上的 Borel 测度. 对于 Borel 集 $E$, 定义 $\mu'(E)=\mu(E^{-1})$, 其中 $E^{-1}=\{g^{-1}\mid g\in E\}$. 证明 $\mu'$ 也是 Borel 测度. 此外, 证明 $\mu$ 是左不变的当且仅当 $\mu'$ 是右不变的.
\end{exercise}

\end{document}